\section{Jenis Kesalahan dan Alat Debugging}
\label{sec:jenis-kesalahan}

Program assembly dapat gagal karena tiga kelompok penyebab. Kelompok
pertama ditemukan sebelum program dijalankan, kelompok kedua baru
terlihat dari hasil yang salah, dan kelompok ketiga membuat program
berhenti sama sekali. Memilih alat dan langkah yang tepat bergantung
pada jenis kesalahannya.

Bagian ini diringkas dari \cite{gunawan2019}, Bagian 4.2--4.3,
hlm.~29--32, dan \cite{sto2001}, BAB XXVII, hlm.~194--202.

\subsection{Tiga Jenis Kesalahan}

\begin{description}
    \item[Kesalahan sintaks] Kesalahan penulisan yang dilaporkan
          assembler, misalnya label belum didefinisikan, ukuran
          operand tidak sama, atau direktif salah tulis. Program tidak
          menghasilkan berkas \file{.OBJ}.

    \item[Kesalahan logika] Program berhasil dirakit dan berjalan,
          tetapi hasilnya tidak sesuai rancangan. Contohnya
          penghitung perulangan dinaikkan padahal seharusnya
          diturunkan, atau nilai dibagi dengan pembagi yang salah.

    \item[Program berhenti (\eng{hang})] Program tidak pernah selesai
          atau melompat ke alamat yang salah. Penyebab yang sering
          terjadi adalah perulangan tak berujung dan stack yang tidak
          seimbang.
\end{description}

\begin{table}[H]
    \centering
    \caption{Jenis kesalahan dan cara menemukannya.}
    \label{tab:jenis-kesalahan}
    \small
    \begin{tabular}{@{}p{3.0cm}p{5.0cm}p{5.0cm}@{}}
        \toprule
        \textbf{Jenis} & \textbf{Gejala} & \textbf{Alat utama} \\
        \midrule
        Sintaks & Pesan kesalahan TASM; tidak ada \file{.OBJ} &
        Daftar kesalahan assembler \\
        \addlinespace
        Logika & Program selesai, hasil salah &
        \eng{Single step}, pemeriksaan register \\
        \addlinespace
        \eng{Hang} & Program tidak kembali ke DOS &
        \eng{Single step}, pemeriksaan stack \\
        \bottomrule
    \end{tabular}
\end{table}

\subsection{Alat Debugging}

\begin{description}
    \item[DEBUG] \eng{Debugger} baris perintah pada DOS. Perintah
          \file{R} menampilkan register, \file{T} menjalankan satu
          instruksi, \file{P} menjalankan satu instruksi dan melewati
          prosedur, \file{D} menampilkan isi memori, dan \file{G}
          menjalankan program sampai alamat tertentu.

    \item[emu8086] Menyediakan tombol \eng{Single Step} serta jendela
          register, \eng{flag}, dan memori. Perubahan nilai dapat
          langsung dilihat setelah setiap instruksi.

    \item[Turbo Debugger] Menyediakan tampilan register, \eng{flag},
          memori, dan stack. Tombol \file{F7} menjalankan satu
          instruksi, sedangkan \file{F8} melewati pemanggilan
          prosedur.
\end{description}

\begin{table}[H]
    \centering
    \caption{Perintah debugging yang sering dipakai.}
    \label{tab:perintah-debug}
    \small
    \begin{tabular}{@{}p{1.8cm}p{5.6cm}p{5.4cm}@{}}
        \toprule
        \textbf{Perintah} & \textbf{Fungsi} & \textbf{Pemakaian} \\
        \midrule
        \file{R} & Menampilkan isi register dan \eng{flag} &
        Memeriksa keadaan sebelum dan sesudah instruksi \\
        \addlinespace
        \file{T} & Menjalankan satu instruksi &
        Masuk ke dalam prosedur \\
        \addlinespace
        \file{P} & Menjalankan satu instruksi, melewati prosedur &
        Menghindari penelusuran prosedur yang sudah benar \\
        \addlinespace
        \file{D} & Menampilkan isi memori &
        Memeriksa variabel, string, dan stack \\
        \addlinespace
        \file{G} & Menjalankan sampai alamat tertentu &
        Melewati bagian yang tidak dicurigai \\
        \bottomrule
    \end{tabular}
\end{table}

\subsection{Langkah Debugging}

Debugging dilakukan secara sistematis, bukan dengan menebak.

\begin{figure}[H]
    \centering
    \begin{tikzpicture}[
        tahap/.style={
            draw=obeBlue, line width=0.8pt, fill=obeBlueLight,
            rounded corners=3pt, minimum width=3.1cm,
            minimum height=1.0cm, align=center, font=\small,
            inner sep=3pt
        },
        panah/.style={-{Stealth[length=2.2mm]}, line width=0.9pt,
            obeTeal}
    ]
        \node[tahap] (rep) {Reproduksi\\kesalahan};
        \node[tahap, right=0.5cm of rep] (iso) {Isolasi\\bagian};
        \node[tahap, right=0.5cm of iso] (hip) {Susun\\hipotesis};
        \node[tahap, right=0.5cm of hip] (uji) {Uji dengan\\\eng{trace}};
        \node[tahap, below=0.8cm of uji] (fix) {Perbaiki};
        \node[tahap, left=0.5cm of fix] (ver) {Verifikasi};

        \draw[panah] (rep) -- (iso);
        \draw[panah] (iso) -- (hip);
        \draw[panah] (hip) -- (uji);
        \draw[panah] (uji) -- (fix);
        \draw[panah] (fix) -- (ver);
    \end{tikzpicture}
    \caption{Alur debugging: reproduksi, isolasi, hipotesis, uji,
    perbaikan, dan verifikasi.}
    \label{fig:alur-debugging}
\end{figure}

\begin{infobox}
\vspace{1.1em}
Satu-satunya bukti bahwa perbaikan berhasil adalah program berjalan
benar dan diuji kembali. Menghapus satu baris yang dicurigai tanpa
memahami penyebabnya hanya memindahkan kesalahan ke tempat lain.
\end{infobox}
